%=============================================================================!
% 14.3  PRESSURE IN A PERIODIC BOX                                            !
%=============================================================================!
\section{Pressure in a periodic box}
\label{sec:pr-periodic}

A periodic box (\cref{sec:md-periodic}) has no walls for its atoms to push
on, and \cref{sec:pr-virial} needed walls. This section finds the
pressure of a periodic box from the free energy, shows that it is
\cref{eq:pr-virial} again with each pair's separation taken to its
nearest copy, and shows what goes wrong when the virial is computed from
the positions instead.

\subsection*{From the free energy}

Chapter~12 found $\dd F = -S\,\dd T - P\,\dd V + \mu\,\dd N$
(\cref{sec:sm-other}), so at fixed $T$ and $N$ the pressure is the rate
at which the free energy falls as the box grows,
\begin{equation*}
   P = -\frac{\partial F}{\partial V} = \kB T\,\frac{\partial\ln
   Z}{\partial V} ,
\end{equation*}
using $F = -\kB T\ln Z$. For a cubic box of side $L$, $V = L^3$, the
partition function is a constant $C$ that depends on $T$ (from the
integral over the momenta) times the integral of $\mathrm{e}^{-\beta U}$
over the positions of every atom inside the box. Writing each position
as $\vb{r}_i = L\vb{s}_i$, with $\vb{s}_i$ the fractional coordinates of
\cref{sec:md-periodic}, each running from 0 to 1, changes the variable of
each of the $3N$ integrals and multiplies it by $L$ (\cref{sec:md-ewald}):
\begin{equation*}
   Z = C\,V^N\int\mathrm{e}^{-\beta U(L\vb{s}_1, \dots, L\vb{s}_N)}\,
   \dd\vb{s}_1\cdots\dd\vb{s}_N .
\end{equation*}
Now the volume appears in two places, $V^N$ and $U$, and the box's
boundaries no longer depend on it. The logarithm is $\ln Z = \ln C +
N\ln V + \ln\int\mathrm{e}^{-\beta U}\,\dd\vb{s}$, with $\dd\vb{s}$ short
for $\dd\vb{s}_1\cdots\dd\vb{s}_N$, and its slope is
\begin{align*}
   \frac{\partial\ln Z}{\partial V}
   &= \frac NV + \frac{\int(-\beta\,\partial U/\partial V)\,
   \mathrm{e}^{-\beta U}\,\dd\vb{s}}{\int\mathrm{e}^{-\beta U}\,\dd\vb{s}}
   && \text{the slope of $\ln$, and the chain rule,}\\
   &= \frac NV - \beta\Bigl\langle\frac{\partial U}{\partial V}\Bigr\rangle
   && \text{a canonical average,}
\end{align*}
with the slope of $U$ taken at fixed fractional coordinates, that is,
with the whole arrangement stretched evenly. For pairs, $U =
\sum_{i<j}\varphi(r_{ij})$, and stretching scales every separation with
the box, $r_{ij} = L\,\lvert\vb{s}_j - \vb{s}_i\rvert$ with the
fractional separation taken to its nearest copy (\cref{sec:md-image}).
So $\partial r_{ij}/\partial L = r_{ij}/L$, and with $\dd L/\dd V =
1/(3L^2)$,
\begin{equation*}
   \frac{\partial U}{\partial V} = \sum_{i<j}\varphi'(r_{ij})\,
   \frac{r_{ij}}{L}\,\frac{1}{3L^2}
   = \frac{1}{3V}\sum_{i<j}r_{ij}\,\varphi'(r_{ij}) = -\frac{\mathcal{W}}{3V} .
\end{equation*}
Then $P = \kB T(N/V - \beta\ens{\partial U/\partial V}) = N\kB T/V +
\ens{\mathcal{W}}/3V$, \cref{eq:pr-virial} again, with $\mathcal{W}$ written
with the separations of \cref{eq:pe-virial}, each to the nearest copy of
the other atom. A periodic run keeps its total momentum at zero, which
removes three of the $3N$ kinetic freedoms, so that $\ens{2K} = (3N -
3)\kB T$ (\cref{sec:sm-equipartition}), and holds its centre of mass
still, so that only $N - 1$ atoms are placed freely and $V^N$ becomes
$V^{N-1}$. The pressure of such a run is $(N - 1)\kB T/V +
\ens{\mathcal{W}}/3V$, one part in $N$ less in its first term, and the
instantaneous $(2K + \mathcal{W})/3V$ averages to it. In
\texttt{md.PairModel} the virial is summed over the pairs of the
neighbour list from their separations, as \cref{sec:pr-tensor} shows,
and \texttt{virial.pressure} adds the kinetic part.

\subsection*{Not from the positions}

The sum $\sum_i\vb{r}_i\cdot\vb{F}_i$ equals \cref{eq:pe-virial} only when,
for every pair, the difference of the positions held, $\vb{r}_j -
\vb{r}_i$, is the separation from which their force was found. Positions
read from the frames of a periodic run are wrapped into the cell
(\cref{sec:md-periodic}): an atom that leaves through one face comes back
through the opposite one, its $\vb{r}_i$ jumping by a lattice vector
while its force does not change. \Cref{sec:pe-forces} showed that moving
the origin by the same vector for every atom leaves $\mathcal{W}$ alone;
wrapping moves single atoms, and changes $\sum_i\vb{r}_i\cdot\vb{F}_i$ by
the scalar product of the lattice vector with that atom's force. Which
atoms have been wrapped depends on where the cell's faces lie, so the sum
depends on where the cell's origin sits. For the liquid of Chapter~13 at
fixed volume, the pressure from the pair separations is
\SI{0.0877}{\giga\pascal}; with $\sum_i\vb{r}_i\cdot\vb{F}_i$ from the
wrapped positions it is \SI{0.0324}{\giga\pascal}, and
\SI{0.0354}{\giga\pascal} with the origin moved, a single frame's value
changing by up to \SI{0.131}{\giga\pascal}. Both lie near the kinetic part
alone, \SI{0.0383}{\giga\pascal}: on average the jumps of the wrapped
atoms cancel the virial's share. Thompson, Plimpton and
Mattson give the general form of the pressure for many-body potentials in
periodic cells, built in the same way from separations\cite{Thompson2009}.

\subsection*{A crystal at rest}

A crystal with every atom on its site has $K = 0$ and its pressure is the
virial's alone. \Cref{fig:pr-crystal}(a) shows it for the face-centred
cubic crystal of the switched Lennard-Jones argon against the cubic
spacing $a$. It is zero at $a_0 = \SI{5.26714}{\angstrom} = 1.54916\sigma$,
the spacing at which \cref{sec:md-start} found the energy least, as it
must be, since $P = -\dd U/\dd V$ at $T = 0$ and the slope of $U$ vanishes
at its least value. A crystal squeezed to $a = \SI{5.20}{\angstrom}$
pushes outwards with \SI{0.128}{\giga\pascal}; one stretched to
\SI{5.40}{\angstrom} pulls inwards with \SI{0.160}{\giga\pascal}, a
negative pressure, a tension. The rise of the pressure for each fraction
by which the volume falls is the bulk modulus,
\begin{equation*}
   B = -V\,\frac{\dd P}{\dd V} = -\frac{\dd P}{\dd V/V} ,
\end{equation*}
the stiffness of the material against squeezing, \SI{2.856}{\giga\pascal}
at $a_0$; its inverse, $\kappa_T = 1/B$ at fixed temperature, is the
compressibility, the fraction by which the volume shrinks for each unit
of pressure.

\begin{figure}[htb]
\centering
\includegraphics{ch14_pressure/crystal}
\caption{The pressure of the face-centred cubic crystal at rest, 256
atoms of the switched Lennard-Jones argon. (a) The pressure (teal) and the
energy per atom (ochre) against the cubic spacing $a$, the spacing $a_0$
of zero pressure marked (dotted). (b) The crystal at $a_0$ stretched along
$x$ by the strain $\epsilon_{xx}$: the components $\mathsf{P}_{xx}$,
$\mathsf{P}_{yy} = \mathsf{P}_{zz}$ and $\mathsf{P}_{xy}$ (dashed, zero) of the
pressure tensor of \cref{sec:pr-tensor}.}
\label{fig:pr-crystal}
\end{figure}

The liquid of Chapter~13, at \SI{135}{\kelvin} and $0.8\sigma^{-3}$, has
the pressure \SI{0.0877}{\giga\pascal}, about 880 bar, of which
\SI{0.0383}{\giga\pascal} is the kinetic part $2K/3V$. Between walls the
same atoms gave \SI{0.0628}{\giga\pascal} (\cref{sec:pr-virial}): the
walls change the liquid next to them, and only the periodic box gives the
pressure of the liquid far from any wall. For the Lennard-Jones argon with
its energy shifted to zero at \SI{8.5}{\angstrom} in place of the switch,
\texttt{mdlab} and ASE's \texttt{LennardJones} give the same pressure,
\SI{-0.058924}{\giga\pascal}, for a strained crystal of 256 atoms, each
displaced at random by about \SI{0.1}{\angstrom} along each axis.

\begin{keyidea}
In a periodic box $P = -\partial F/\partial V$ gives \cref{eq:pr-virial}
again, with the virial summed over pairs from their separations to the
nearest copy; $\sum_i\vb{r}_i\cdot\vb{F}_i$ from wrapped positions
changes with the origin and is wrong. A crystal at rest has zero
pressure at the spacing of least energy, and its bulk modulus $B =
-V\,\dd P/\dd V$ measures its stiffness against squeezing.
\end{keyidea}

\notebook{14}{shift the origin of a periodic box and watch which pressure changes}

\begin{selfcheck}
From \cref{fig:pr-crystal}(a), the pressure falls from +0.128 to
\SI{-0.160}{\giga\pascal} as $a$ goes from 5.20 to \SI{5.40}{\angstrom}.
Estimate the bulk modulus near $a_0$ from these two points.
\end{selfcheck}
